% =====================================================================
%  Wspólna preambuła szablonów pism ZUS (świadczenia rodzicielskie 2026)
%  Silnik: XeLaTeX albo Tectonic (pdflatex NIE obsłuży polskich znaków w fontspec).
%  Kompilacja:  tectonic -X compile plik.tex
% =====================================================================
\documentclass[11pt,a4paper]{article}

\usepackage{fontspec}
\usepackage{polyglossia}
\setmainlanguage{polish}
\setotherlanguage{english}
% font domyślny (Latin Modern) — Tectonic go dostarcza

\usepackage[margin=2.5cm,a4paper]{geometry}
\usepackage{graphicx}
\usepackage{xcolor}
\usepackage{array}
\usepackage{enumitem}
\usepackage{fancyhdr}
\usepackage{hyperref}

\definecolor{zus}{HTML}{007A33}      % zieleń ZUS
\definecolor{szary}{HTML}{6B7280}
\definecolor{ramka}{HTML}{9CA3AF}

\hypersetup{colorlinks=true, linkcolor=zus, urlcolor=zus, pdftitle={Szablony pism ZUS}}

% --- pole do wypełnienia odręcznego -----------------------------------
% \Pole{nazwa} rysuje kropkowaną linię z podpisem pola pod spodem.
% Pole do wypełnienia: kropkowana linia, pod nią drobny podpis pola.
\newcommand{\Pole}[1]{%
  \textcolor{szary}{\dotfill}\,%
  {\scriptsize\textcolor{szary}{\textsc{#1}}}%
}

% --- blok warunkowy: treść pokazywana tylko, gdy pole prawdziwe --------
\newenvironment{warunek}{%
  \par\smallskip\noindent\begin{minipage}{\linewidth}%
  \color{black}%
}{%
  \end{minipage}\par\smallskip%
}

% --- nagłówek pisma ---------------------------------------------------
\newcommand{\NaglowekPisma}[1]{%
  \par\vspace{0.8em}%
  \begin{center}%
    {\large\bfseries\MakeUppercase{#1}}%
  \end{center}%
  \vspace{0.4em}%
}

% --- blok nadawcy (pierwsze linie pisma) -------------------------------
\newenvironment{nadawca}{%
  \begin{flushleft}\small
}{%
  \end{flushleft}\normalsize
}

% --- stopka z numerem strony -----------------------------------------
\pagestyle{fancy}
\fancyhf{}
\renewcommand{\headrulewidth}{0pt}
\fancyfoot[C]{\textcolor{szary}{\small\thepage}}

% --- metadane szablonu na początku pisma ------------------------------
\newcommand{\Metadane}[3]{%
  \begin{flushright}\footnotesize\textcolor{szary}%
    \textbf{#1}\\[2pt] {\itshape #2}\\[2pt] #3%
  \end{flushright}\vspace{0.5em}%
}

\setlength{\parindent}{0pt}
\setlength{\parskip}{0.6em}

\begin{document}

\begin{flushright}
{\footnotesize\color{szary}
\textbf{Szablon P12}\\
{\itshape Wniosek o urlop na warunkach urlopu macierzyńskiego (przysposobienie)}\\
art. 183 \S{} 1 k.p.; art. 29a, art. 30 ustawy zasiłkowej -- Formularz: ZAM}\
\end{flushright}

\vspace{0.5em}
{\color{ramka}\hrule height 0.4pt}
\vspace{0.8em}

\NaglowekPisma{WNIOSEK O UDZIELENIE URLOPU NA WARUNKACH URLOPU MACIERZYŃSKIEGO}

Na podstawie art. 183 \S{} 1 Kodeksu pracy wnoszę o udzielenie mi urlopu na
warunkach urlopu macierzyńskiego w wymiarze \Pole{tygodnie} tygodni, w okresie
od \Pole{urlop od} do \Pole{urlop do}, w związku z:

[ ] przyjęciem dziecka na wychowanie w celu przysposobienia,
[ ] przyjęciem dziecka na wychowanie w ramach rodziny zastępczej
(innej niż zawodowa).

Dziecko: \Pole{dziecko imie}, urodzone \Pole{data porodu}.
Przyjęcie dziecka na wychowanie nastąpiło dnia \Pole{data zdarzenia}.

\begin{warunek}
Urlop przysługuje mi w wymiarze 20 tygodni (przy jednym dziecku), a przy\\[-2pt]
przyjęciu dziecka w wieku do 3 lat --- do ukończenia przez dziecko 4 lat.
\end{warunek}

\begin{warunek}
Urlop przysługuje mi w wymiarze wynikającym z art. 183 \S{} 1 k.p.;\\[-2pt]
rodzina zastępcza zawodowa nie ma prawa do tego urlopu.
\end{warunek}

\Pole{data pisma}
.....................................
(podpis)

\vfill
{\footnotesize\textcolor{szary}\itshape Adopcja i rodzina zastępcza mają RÓŻNE wymiary urlopu. \quotedblbase{}Za życiem\textquotedblright{} (art. 183 \S{} 4 pkt 3 i \S{} 4\textsuperscript{1} k.p.) dotyczy rodziny zastępczej, nie przysposobienia. }

\end{document}
